\documentclass[11pt]{article}
\usepackage{preamble}

\begin{document}

\packetheading{Lesson 4-3 \textperiodcentered\ Period 2 \textperiodcentered\ Student Packet}{Multiply \& Divide Rational Expressions + DOK-3 Closure}

\sectionbanner{OBJECTIVES}
{\small
\renewcommand{\arraystretch}{1.25}
\noindent\begin{tabular}{|p{1.8in}|p{4.8in}|}
\hline
\textbf{Math Objective} & Multiply and divide rational expressions by factoring, canceling common factors, and rewriting division as multiplication by the reciprocal; track all domain restrictions through every step. \\ \hline
\textbf{Language Objective} & Justify in writing whether a set is closed under an operation, using sentence frames that cite factoring and cancellation. \\ \hline
\textbf{Essential Understanding} & Rational expressions behave like rational numbers: closed under multiplication and nonzero division, because factoring + cancellation always produces another rational expression. \\ \hline
\end{tabular}
}

\sectionbanner{TOPIC GOALS / ESSENTIAL QUESTION / MATERIALS}
{\small
\renewcommand{\arraystretch}{1.25}
\noindent\begin{tabular}{|p{1.8in}|p{4.8in}|}
\hline
\textbf{Topic Goals} & Multiply and divide rational expressions; justify closure under these operations. \\ \hline
\textbf{Essential Question} & Why are rational expressions closed under multiplication and division, and what must we track to prove it? \\ \hline
\textbf{Materials} & Student packet, pencil. Teacher models Example 3; DOK-3 Practice \#13 is the capstone closure-argument task. \\ \hline
\end{tabular}
}

\sectionbanner{LAUNCH --- Multiply \& Divide Rational Expressions (teacher models Example 3)}

\begin{callout}{\IconBook\ MULTIPLY / DIVIDE RULES (keep printed on your page)}{calloutgreen}
MULTIPLY / DIVIDE RULES
\begin{enumerate}[leftmargin=2.1em,itemsep=1pt,topsep=4pt]
  \item Factor every numerator AND denominator completely.
  \item For division, rewrite as multiplication by the RECIPROCAL of the divisor.
  \item Cancel only common FACTORS across any numerator with any denominator.
  \item State the domain: exclude every zero of every denominator that appeared in ANY step, including the reciprocal's new denominator.
\end{enumerate}
\end{callout}

\begin{callout}{\IconWarn\ DIVISION IS MULTIPLICATION BY THE RECIPROCAL}{calloutred}
\[
\frac{a}{b}\div \frac{c}{d} = \frac{a}{b}\times \frac{d}{c}
\]
The new denominator (\(c\)) must be added to your domain-restriction list even if it cancelled.\\
Never cancel across the \(\div\) sign --- flip FIRST, then cancel.
\end{callout}

\sectionbanner{EXPLORE --- Think-Pair-Share (33 min)}
\textit{Work with a partner. Move in order. Practice \#13 (Closure Argument) is the big one --- plan \(\sim\)12 min for it.}

\textbf{Bank-item labels (in order):}
\begin{enumerate}[leftmargin=1.6em,itemsep=2pt,topsep=2pt]
  \item Try It 3 --- Multiply rational expressions
  \item Practice \#27 --- Multiply with factoring
  \item Try It 4 --- Divide rational expressions
  \item Practice \#29 --- Divide: rewrite as multiplication
  \item Try It 5 --- Mixed multiply/divide
  \item Practice \#32 --- Multi-step multiply/divide
  \item Practice \#13  \IconStar\ DOK-3 CLOSURE ARGUMENT
\end{enumerate}

\bankitem{Try It 3 --- Multiply rational expressions}{%
Find the simplified form of each product, and state the domain.

\begin{enumerate}[label=\textbf{\alph*.},leftmargin=1.8em,itemsep=3pt,topsep=4pt]
  \item \(\displaystyle \frac{x^2-16}{9-x}\cdot \frac{x^2+x-90}{x^2+14x+40}\)
  \item \(\displaystyle \frac{x+3}{4x}\cdot \frac{3x-18}{6x+18}\cdot \frac{x^2}{4x+12}\)
\end{enumerate}
}{}

\bankitem{Practice \#27 --- Multiply with factoring}{%
Find the product and the domain.

\[
\frac{(x-y)^2}{x+y}\cdot \frac{3x+3y}{x^2-y^2}
\]
}{}

\bankitem{Try It 4 --- Divide rational expressions}{%
Find the simplified form of each product and the domain.

\begin{enumerate}[label=\textbf{\alph*.},leftmargin=1.8em,itemsep=3pt,topsep=4pt]
  \item \(\displaystyle \frac{x^3-4x}{6x^2-13x-5}\cdot (2x^3-3x^2-5x)\)
  \item \(\displaystyle \frac{3x^2+6x}{x^2-49}\cdot (x^2+9x+14)\)
\end{enumerate}
}{}

\bankitem{Practice \#29 --- Divide: rewrite as multiplication}{%
Find the product and the domain.

\[
\frac{2x^2-10x}{(x-5)(x^2-1)}\cdot (3x^2+4x+1)
\]
}{}

\bankitem{Try It 5 --- Mixed multiply/divide}{%
Find the simplified quotient and the domain of each expression.

\begin{enumerate}[label=\textbf{\alph*.},leftmargin=1.8em,itemsep=3pt,topsep=4pt]
  \item \(\displaystyle \frac{1}{x^2+9x}\div \left(\frac{6-x}{3x^2-18x}\right)\)
  \item \(\displaystyle \frac{2x^2-12x}{x+5}\div \left(\frac{x-6}{x+5}\right)\)
\end{enumerate}
}{}

\bankitem{Practice \#32 --- Multi-step multiply/divide}{%
Find the quotient and the domain.

\[
\frac{25x^2-4}{x^2-9}\div \frac{5x-2}{x+3}
\]
}{}

\bankitem{Practice \#13  \IconStar\ DOK-3 CLOSURE ARGUMENT}{%
\textbf{Higher Order Thinking} Why are rational expressions closed under multiplication and division?

\vspace{1.4in}
}{}

\sectionbanner{SHARE / SUMMARY (5 min)}

\begin{sentenceframebox}
\textbf{Sentence frames}

Pair share (Closure Argument): ``I rewrite the division as multiplication by the reciprocal, giving \_\_\_. After factoring, I cancel the factor \_\_\_, leaving \_\_\_. The domain excludes \(x=\) \_\_\_ because those values make at least one denominator zero at some step.''

Whole class: one pair reads the closure argument. Class signals thumbs up/sideways.
\end{sentenceframebox}

\begin{callout}{\IconSearch\ BRIDGE TO P3}{calloutpeach}
Next period we apply multiply/divide skills to complex rational expressions and assessment-critical items. Bring your domain-restriction list habit.
\end{callout}

\sectionbanner{EXIT}

\begin{summaryexitbox}{EXIT TICKET --- Summary Recap}
To divide rational expressions I first \_\_\_.

I cancel common factors, not \_\_\_.

My domain restriction list must include zeros from \_\_\_.

One biggest thing I learned today: \_\_\_\_\_\_\_\_\_\_\_\_\_\_\_\_\_\_\_\_\_\_\_\_\_\_
\end{summaryexitbox}

\clearpage

\sectionbanner{OPTIONAL REINFORCEMENT (not collected)}
\textit{If you finish Explore early. \#40 is a finance performance task that applies rational expressions to a real-world context.}

\bankitem{Optional \#1  (Savvas Practice, performance-task finance item)}{%
\textbf{Performance Task} The approximate annual interest rate \(r\) of a monthly installment loan is given by the formula:

\[
r=\frac{\frac{24(nm-p)}{n}}{p+\frac{nm}{12}}
\]

where \(n\) is the total number of payments, \(m\) is the monthly payment, and \(p\) is the amount financed.

Part A Find the approximate annual interest rate (to the nearest percent) for a four-year signature loan of \$20{,}000 that has monthly payments of \$500.

Part B Find the approximate annual interest rate (to the nearest tenth percent) for a five-year auto loan of \$40{,}000 that has monthly payments of \$750.
}{}

\end{document}
